\section{Experiments: Does the Gap Appear Beyond Witnesses?}
\label{sec:experiments}

The theoretical separations in
\Cref{sec:minimal-1nn,sec:odd-k-separation,sec:even-k-separation} are
worst-case constructions.  This section asks: how often do LOO/replace-one gaps
appear in finite samples that were not hand-crafted?  The experiments are not
intended to prove a theorem; they provide reproducible evidence that the
separation mechanism is detectable, not confined to the minimal witnesses.

\subsection{Metrics}

We use two primary metrics.

\begin{description}[leftmargin=*,style=nextline]
    \item[Strict separation frequency]
        $\Pr[\max_i \Delta_{\mathrm{loo}}(S,i) = 0 \;\text{and}\;
        \max_{j,z,x,y} \Delta_{\mathrm{rep}}(S,j,z,x,y) = 1]$.
        The proportion of random samples for which LOO is perfectly stable
        while a replace-one perturbation changes the loss at some query.
    \item[Vulnerable-query rate]
        $\Pr_x[|\M_k(S,x)| \le 2]$.
        The fraction of metric-space points whose signed vote margin lies
        within $\pm 2$, making them potentially vulnerable to a single
        label-flip replacement (\Cref{prop:5.1}).  A high vulnerable-query rate
        indicates that the separation mechanism is not pathologically rare.
\end{description}

All experiments fix deterministic tie-breaking (distance, then sample index;
label ties resolved $0 \prec 1$) and binary 0--1 loss.

\subsection{Synthetic graph metrics}

We generate random connected graphs on 2--5 vertices via the
Erd\H{o}s--R\'enyi model with edge probability $0.5$, convert them to
shortest-path metric spaces, and draw labeled samples with controlled duplicate
ratio, conflict ratio, and label noise.  For each configuration we run 50
independent trials.

On 2-vertex graphs with zero noise and zero conflict, the separation frequency
decreases as $k$ grows and increases with the duplicate ratio, consistent with
the theoretical mechanism: duplicates create more opportunities for
margin-zero configurations where a single label flip crosses the decision
boundary.  At $k=1$ with duplicate ratio $0.6$, the separation reaches $0.32$,
and it rises to $0.48$ at $k=5$.

\subsection{UCI tabular data}

We evaluate the same metrics on four UCI datasets: Iris, Wine, Breast Cancer
Wisconsin, and Digits.  Each dataset is converted to binary classification
(class~0 vs.\ the rest) and embedded in a Euclidean metric space.  For each
dataset, we draw 50 random subsamples of up to 100 points and compute the
separation frequency and vulnerable-query rate.

At $k=1$, Iris and Digits exhibit a separation frequency of $1.0$---every
subsample has a replace-one flip while all LOO indicators stay zero.  Wine and
Breast Cancer show no separation at $k=1$ (separation $=0.0$), because LOO is
already unstable on those datasets.  The vulnerable-query rate is $1.0$ at
$k=1$ for all datasets, as every query has $|\margin|=1$ under 1-NN.  At higher
$k$, the separation frequency drops sharply ($\le 0.14$ for $k=3,5$),
confirming that the gap is concentrated at the smallest neighborhood size
where the margin is most fragile.

\subsection{Margin-based vulnerability diagnostic}

The margin-based diagnostic ($|\M_k(S,x)| \le 2$ flags a query as potentially
vulnerable) is not merely a heuristic: for odd $k$, a single replacement
alters the ordered vote by at most $\pm 2$, so the margin can cross zero if and
only if $|\M_k(S,x)| \le 2$.  This follows directly from the fact that
replacing one occurrence changes the top-$k$ set by at most two entries, each
contributing $\pm 1$ to the signed margin.  Consequently the diagnostic has
precision and recall identically $1$ on all odd-$k$ configurations---no
empirical validation is needed.  For even $k$ the situation is more complex
because the margin is always even and the $0 \prec 1$ tie-break creates
additional asymmetry.

The practical upshot is that practitioners can efficiently identify potentially
vulnerable queries by computing the margin distribution rather than relying
solely on aggregate LOO error.

\subsection{Empirical summary}

\begin{enumerate}[leftmargin=*,label=(\arabic*)]
    \item The LOO/replace-one separation appears at measurable frequency in
        random graph metrics under realistic duplicate ratios.
    \item The gap is not confined to the minimal two-point witness: on Digits
        at $k=1$, essentially every subsample exhibits the separation.
    \item The vulnerable-query rate ($|\M_k| \le 2$) is a reliable diagnostic
        for odd $k$.
    \item Duplicates are the primary structural driver; noise and conflict play
        secondary roles.
\end{enumerate}

These findings support the claim that the LOO/replace-one separation is a
structurally stable phenomenon detectable beyond hand-crafted examples.  Full
experimental protocol, parameters, and random seeds are reported in
\Cref{app:experimental-protocol}.
